\chapter{Interfaces that preserve experimental meaning}
\label{completion:composition}

A local equivalence certificate is useful in a larger apparatus only when
the allowed coupling respects what that certificate preserves. This chapter
states a finite classical interface contract, proves its error propagation,
and gives an exact refinement algorithm. The ingredients are established
stochastic lumpability, intertwining and partition refinement. Their role
here is to make modular use of access certificates checkable.

\section{Joint outcome and interface state}
A finite classical instrument has nonnegative matrices $M_{a,y}$ with
$\sum_yM_{a,y}$ stochastic. Let $C$ be an $n\times k$ row-stochastic
interface, and $B_{a,y}$ a $k$-state instrument with the same actions and
outcomes. Define the local defect
\[
 \epsilon_a(x)=\frac12\sum_{y,z}
       |(M_{a,y}C-CB_{a,y})_{xz}|.
\]
The defect includes both the recorded outcome and the state passed to later
operations. Agreement of immediate readout probabilities alone is not enough.
Preparations, couplings, controls and readout updates belong to the contract.

\begin{theorem}[Adaptive composition bound]
Let $P,Q$ be the complete terminal-record laws of the concrete and quotient
instruments under the same finite adaptive policy tree, with respective
initial rows $\mu,\nu$. At each internal history $h$, let
$\mu_h=\mu M_h$ be the unnormalized concrete branch row and $a(h)$ its
chosen action. Then
\[
 \operatorname{TV}(P,Q)\leq
 \min\left\{1,\operatorname{TV}(\mu C,\nu)
                  +\sum_h\mu_h\epsilon_{a(h)}\right\}.
\]
In particular, exact intertwining $M_{a,y}C=CB_{a,y}$ and matching interface
preparations imply identical records for every permitted finite adaptive
policy, including outcome-dependent stopping.
\end{theorem}
\begin{proof}
For a quotient branch row $\nu_h$, put $d_h=\mu_hC-\nu_h$. At an action $a$,
\[
 d_{hy}=\mu_h(M_{a,y}C-CB_{a,y})+d_hB_{a,y}.
\]
Nonnegativity gives
$\sum_y\|\mu_h(M_{a,y}C-CB_{a,y})\|_1\leq2\mu_h\epsilon_a$.
Outcome normalization gives
$\sum_y\|d_hB_{a,y}\|_1\leq\|d_h\|_1$.
Sum through the tree, retaining every outcome and stopped branch. The leaf
probability difference is $d_h\mathbf1$ because $C\mathbf1=\mathbf1$.
Half the sum of absolute leaf probability differences proves the claim.
\end{proof}

The defect sum is an expected accumulated error under the concrete model;
it can improve on horizon times a uniform worst-case defect. Exact
intertwining is a sufficient contextual certificate. It is not asserted to
characterize every trace-equivalent pair.

\section{Independent components and new couplings}
For product instruments with product interface $C_1\otimes C_2$, the local
joint defect at a state pair obeys
\[
 \epsilon_{12}\leq\epsilon_1+\epsilon_2-\epsilon_1\epsilon_2.
\]
Maximally couple each pair of normalized outcome/interface distributions
independently. Both agree with probability at least
$(1-\epsilon_1)(1-\epsilon_2)$, proving the bound. For several factors use
$1-\prod_i(1-\epsilon_i)$. A subsequent adaptive context is covered by the
preceding theorem if every added interaction also has a checked joint defect.

Matching local preparation marginals does not establish matching joint
preparations. The distributions $(00+11)/2$ and $(01+10)/2$ have identical
one-bit marginals; a controlled-NOT followed by target-bit readout identifies
their parity perfectly. Joint interface equality must be established,
for example through an independently justified product preparation.

\section{Constructing a sufficient deterministic interface}
Given an initial observation partition $\mathcal P_0$, repeatedly split each
block according to the signature
\[
 \left(\sum_{j\in B}M_{a,y}(i,j)\right)_{a,y,B\in\mathcal P}.
\]
\begin{proposition}
The iteration terminates after at most $n-|\mathcal P_0|$ strict rounds and
returns the coarsest deterministic Markov interface refining $\mathcal P_0$
for all declared instruments and all initial preparations.
\end{proposition}
\begin{proof}
Each strict round increases the number of blocks. At termination, states in
one block send the same outcome/transition mass to every target block, exactly
the required intertwining identity for the deterministic partition matrix.
Any stable refinement of the current partition must separate unequal
signatures, since each current block is a union of refined blocks. Induction
shows that every stable refinement refines every iteration, and hence the
terminal partition. This proves coarseness as well as closure.
\end{proof}

The software exports the forced refinements and independently checks every
split and final closure. Minimality refers to deterministic Markov interfaces;
it does not imply a smallest linear realization or a cheapest physical sensor.
An algorithm naming a finer hidden block does not supply physical access to it.

In the four-state model $(h,o)$, passive readout sees only $o$. Preparations
$(0,0)$ and $(1,0)$ are passively identical. Adding the allowed reversible
logical coupling $(h,o)\mapsto(h,o\mathbin{\mathrm{XOR}}h)$ makes a subsequent
visible readout distinguish them. The coarsest interface expands from two
states to four. This illustrates an actual change in permitted access, while
making no claim that physical implementation of the logical gate costs zero
energy. A second example has identical immediate output laws but different
hidden updates; its later records separate. The joint defect detects that
failure of a readout-only calibration.

The exact implementations are in
\path{constructive_access/completion/composition.py}; reproducible examples
and independent checks accompany the source. No theorem in this chapter is
claimed for arbitrary quantum couplings or unmodeled shared environments.
